% ==============================================================================
% T0/FFGFT Theory: Document 222
% Title: Lithium falsification criteria:
%        FFGFT predictions for light element abundances
% Author: Johann Pascher
% Date: May 2026
% References: Documents 025 (Cosmology), 063 (Cosmic), 220 (Casimir falsification),
%             221 (Redshift falsification)
% Occasion: P. Austin (Information Physics Institute), falsification request
% ==============================================================================

\section*{Preamble}

This document formulates the falsification criterion for the FFGFT
position on the cosmological lithium problem, in direct response to
P.~Austin's methodological question:
\emph{\enquote{What abundance pattern would distinguish the FFGFT
equilibrium interpretation from BBN plus stellar depletion?}}

Lithium is the most difficult of the three falsification tests.
Casimir (Doc.~220) and redshift (Doc.~221) build on local
$\xi_0$-vacuum action; lithium tests the \emph{cosmic prehistory} of
the universe. The FFGFT answer here is not a modification of a single
reaction rate, but a fundamentally different view on the origin of
the light elements.

% ==============================================================================
\section{The cosmological lithium problem}
% ==============================================================================

\subsection{Observational situation}

The principal light-element abundances in metal-poor population II
stars (Spite plateau, $\text{[Fe/H]} < -2$):

\begin{center}
\renewcommand{\arraystretch}{1.3}
{\small
\begin{tabular}{>{\raggedright\arraybackslash}p{1.6cm}>{\raggedright\arraybackslash}p{2.3cm}>{\raggedright\arraybackslash}p{2.3cm}>{\raggedright\arraybackslash}p{1.8cm}}
\toprule
\textbf{Element} & \textbf{Observed} & \textbf{BBN prediction} & \textbf{Status} \\
\midrule
$^4$He   & $24.8\,\%$ mass fraction    & $24.7$--$24.8\,\%$  & \checkmark{} \\
D/H      & $(2.5 \pm 0.2)\!\cdot\!10^{-5}$  & $(2.5 \pm 0.1)\!\cdot\!10^{-5}$ & \checkmark{} \\
$^7$Li/H & $(1.6 \pm 0.3)\!\cdot\!10^{-10}$ & $(5.0 \pm 0.7)\!\cdot\!10^{-10}$ & \textbf{factor 3} \\
$^6$Li/H & $\lesssim 10^{-12}$ & $\lesssim 10^{-14}$ & marginal \\
\bottomrule
\end{tabular}
}
\renewcommand{\arraystretch}{1.0}
\end{center}

While $^4$He and D/H are consistent with BBN+Planck, $^7$Li is short
by a factor of $3$--$4$. This discrepancy has been unresolved for two
decades.

\subsection{Three standard solution attempts}

\begin{enumerate}[leftmargin=2em,itemsep=2pt,topsep=2pt]
  \item \textbf{Modified reaction rates:} CERN n\_TOF experiments on
        $^7$Be destruction rates. So far no sufficient correction
        found.
  \item \textbf{Stellar depletion:} mixing, diffusion, atmospheric
        depletion in Spite stars. Consistent with $^7$Li loss, but
        the Spite plateau is too smooth for depletion models.
  \item \textbf{Structure formation effects:} Pop-III stellar capture
        (Aanda 2025, Yeh \& Olive 2022). Explains part of the
        discrepancy, depends on Pop-III IMF.
\end{enumerate}

% ==============================================================================
\section{FFGFT position: no BBN phase}
% ==============================================================================

\subsection{Static universe vs.\ standard BBN}

Standard BBN presupposes a \emph{cooling, expanding phase} from
$T \sim 10^{11}$ K to $T \sim 10^{9}$ K in $\sim 3$ minutes. In this
phase:
\begin{itemize}[leftmargin=2em,itemsep=2pt,topsep=2pt]
  \item the deuterium bottleneck at $T \sim 10^{9}$ K is the
        rate-limiting step,
  \item the destruction of $^7$Be by electron capture
        ($t_{1/2} \approx 53$ days) is too slow to complete before
        \emph{cosmological cooldown},
  \item the resulting $^7$Li level is a \emph{frozen-in remnant} of
        an interrupted reaction sequence.
\end{itemize}

In FFGFT this phase does not exist. The universe is static; there is
no cosmological cooldown from $10^{11}$ to $10^{9}$ K in 3 minutes.
The cosmic microwave background has always been at $T_\text{CMB} =
2.725$ K (see Doc.~091, $\rho_\text{CMB} = \xi_0 \hbar c / L_\xi^4$).
At this temperature all atomic nuclei are stable, and there are no
BBN reactions.

\textbf{Consequence:} the observed abundances of $^4$He, D, $^7$Li,
$^6$Li in metal-poor stars are in FFGFT \emph{not} primordial in the
standard sense. They result from a \emph{stellar prehistory} with
unbounded time.

\subsection{Stellar sources of light elements}

In the FFGFT view, the abundances arise from the long-term
equilibrium between astrophysical production and destruction
mechanisms:

\begin{itemize}[leftmargin=2em,itemsep=2pt,topsep=2pt]
  \item \textbf{$^4$He:} primary producer is hydrogen fusion in stars
        of all masses. In a universe with unbounded cosmic time,
        Pop-III stars and successive generations enrich the
        interstellar medium to $\sim 25\,\%$ He. Consistent with
        present observation.
  \item \textbf{D:} D is \emph{consumed} by stellar fusion and
        \emph{produced} by cosmic spallation on helium nuclei.
        Equilibrium gives $D/H \sim 10^{-5}$ -- consistent with
        observation.
  \item \textbf{$^7$Li:} produced in AGB stars, classical novae,
        cosmic spallation. Destroyed in stellar surfaces by
        p-capture. Equilibrium level $\sim 10^{-10}$ -- consistent
        with the Spite plateau.
  \item \textbf{$^6$Li:} almost exclusively from cosmic spallation;
        no BBN contribution expected. Level $\lesssim 10^{-12}$
        consistent with sparse upper bounds.
\end{itemize}

\textbf{What this changes:} the \enquote{lithium problem} disappears
in this view, because there is no BBN prediction value exceeding the
observation. Instead, the task is to show that stellar sources
quantitatively reproduce the Spite plateau level -- which is exactly
what the Aanda 2025 paper and Pop-III modelling are pursuing, albeit
within the expansion paradigm.

% ==============================================================================
\section{FFGFT predictions}
% ==============================================================================

\subsection{\texorpdfstring{Constancy of $\xi_0$ over cosmic time}{Constancy of xi_0 over cosmic time}}

The \emph{most important} FFGFT prediction in the lithium domain is
that $\xi_0 = 4/30000$ \textbf{is a time-independent geometric
constant}. Hence:
\begin{itemize}[leftmargin=2em,itemsep=2pt,topsep=2pt]
  \item the fine-structure constant $\alpha = \xi_0 \cdot E_0^2$
        is constant (see Doc.~202, §22),
  \item all nuclear binding energies (Coulomb part $\propto \alpha$,
        nuclear interaction scale-free) are constant,
  \item all reaction rates (cosmic spallation, stellar
        nucleosynthesis) are constant.
\end{itemize}

\textbf{Falsification:} variation tests of natural constants
(quasar spectrum analyses across large $z$, Webb \& Murphy et al.,
many-multiplet method) show no variation in most studies; an
unambiguous detection $|\Delta\alpha/\alpha| > 10^{-5}$ between
$z = 0$ and $z = 4$ would refute FFGFT. \emph{Currently:} maximum
$|\Delta\alpha/\alpha| < 10^{-6}$, FFGFT consistent.

\subsection{Constancy of element abundances in old structures}

In standard BBN+expansion, all abundances are $\propto \eta$
(baryon-to-photon ratio) and thus time-dependent in the early phase.
In FFGFT, $\eta$ is a position-dependent quantity of the present
cosmic matter distribution, not a time-evolving parameter.

\textbf{FFGFT prediction:} element abundances in old, metal-poor
stars show a \emph{smooth plateau} without systematic dependence on
the redshift of observation. In the standard BBN+expansion picture,
slight variations would be expected at very high redshifts ($z > 6$),
since these observations probe closer to the actual nucleosynthesis.

\textbf{Falsification:} if JWST observations of galaxies at $z > 6$
systematically find \emph{different} abundances than the local Spite
plateau, the FFGFT view (abundance as stellar equilibrium, not
primordial value) is challenged.

\subsection{\texorpdfstring{Quantitative prediction: no $^7$Li discrepancy}{Quantitative prediction: no 7Li discrepancy}}

In FFGFT the factor of $3$ between BBN prediction and observation
\emph{vanishes} -- not through a modified reaction rate, but because
the BBN prediction is not the correct reference. The correct
reference is the \emph{stellar equilibrium}, lying at $^7$Li/H $\sim
10^{-10}$.

\textbf{Falsification:} if the standard BBN prediction can be
reduced to the observed value by corrections to nuclear reaction
rates (e.g.\ new $^7$Be$(n,\alpha)^4$He resonances from n\_TOF)
\emph{without} worsening D/H and $^4$He, the FFGFT position becomes
superfluous (but not refuted). A genuine refutation would be a
\emph{complete, consistent} standard BBN solution for all four
abundances without additional assumptions.

\subsection{\texorpdfstring{$^4$He abundance as critical test}{4He abundance as critical test}}

The \emph{hardest} internal test of the FFGFT position is the $^4$He
abundance of $24.8\,\%$ in extremely metal-poor stars. Standard BBN
explains this number in the first 3 minutes of cosmic history. FFGFT
must explain this number through stellar prehistory.

\textbf{FFGFT requirement:} Pop-III stars (or comparable early
populations) must produce sufficient He enrichment to reach $\sim
25\,\%$ mass fraction, without destroying $^7$Li, D, or metal
abundances in Spite stars. The available cosmic time in FFGFT is
\emph{unbounded}, which opens this path -- but it must be
quantitatively consistent.

\textbf{Falsification:} a Pop-III/II stellar model that
\emph{provably} delivers either too little or too much $^4$He while
simultaneously reproducing the other three abundances.

% ==============================================================================
\section{Specific falsification criteria}
% ==============================================================================

FFGFT is \textbf{refuted} by the following observations:

\begin{enumerate}[leftmargin=2em,itemsep=4pt,topsep=2pt]
  \item \textbf{Variation of $\alpha$ or $\xi_0$:}
        a significant variation $|\Delta\alpha/\alpha| > 10^{-5}$
        between $z = 0$ and high redshifts.
        \emph{Current status:} no variation detected, FFGFT
        consistent.

  \item \textbf{Abundance variation with $z$:}
        if JWST spectroscopy of light elements in galaxies at $z > 6$
        systematically finds different abundances than the local
        Spite plateau, not explainable by local chemical evolution.
        \emph{Current status:} undecided, critical test in coming
        years.

  \item \textbf{Complete standard BBN solution:}
        a consistent modification of nuclear reaction rates (e.g.\
        from n\_TOF) that reproduces all four abundances ($^4$He, D,
        $^7$Li, $^6$Li) \emph{without} additional free parameters
        within the standard expansion framework. This would not
        refute FFGFT geometry, but would render the FFGFT position
        on the lithium problem superfluous.

  \item \textbf{Inconsistency of stellar sources:}
        a quantitative stellar model showing that observed
        abundances cannot be reproduced from stellar sources +
        cosmic spallation without invoking BBN.
        \emph{Current status:} Aanda 2025 work shows that Pop-III
        modelling can explain the Spite plateau. Full consistency
        for all four abundances without BBN is not yet demonstrated.

  \item \textbf{$^4$He variation in old structures:}
        if $^4$He abundance in extremely metal-poor high-redshift
        HII regions systematically deviates from $24.8\,\%$.
        \emph{Current status:} current measurements at moderate
        redshifts all give $\sim 24.8\,\%$.
\end{enumerate}

% ==============================================================================
\section{What FFGFT does NOT predict}
% ==============================================================================

\begin{itemize}[leftmargin=2em,itemsep=2pt,topsep=2pt]
  \item FFGFT does \emph{not} predict a specific numerical value for
        $^7$Li/H = $X$ competing with BBN. The position is
        methodological: BBN is not the reference, because there is
        no BBN phase.
  \item FFGFT does \emph{not} predict any new exotic reaction
        converting $^7$Li. The solution lies in the \emph{history},
        not in the \emph{physics of a single reaction}.
  \item FFGFT does \emph{not} predict any anomaly of stellar fusion
        or spallation rates. Both processes proceed with standard
        cross sections.
\end{itemize}

% ==============================================================================
\section{Connection to Doc.~220 and 221}
% ==============================================================================

The three falsification documents form a hierarchical structure:

\begin{itemize}[leftmargin=2em,itemsep=2pt,topsep=2pt]
  \item \textbf{Doc.~220 (Casimir):} tests $\xi_0$-vacuum on
        $\sim 100\,\mu$m scale. Local, directly measurable.
  \item \textbf{Doc.~221 (Redshift):} tests the same $\xi_0$
        mechanism on cosmological scales. Astronomical, indirect
        but precise.
  \item \textbf{Doc.~222 (Lithium):} tests the \emph{cosmic history}
        -- that the universe is static and $\xi_0$ is constant over
        cosmic time. Model-based, since direct temporal observation
        of element synthesis is impossible.
\end{itemize}

\textbf{Hierarchy of strength:}
\begin{itemize}[leftmargin=2em,itemsep=2pt,topsep=2pt]
  \item Casimir falsification (Doc.~220) is the \emph{strongest} --
        immediately measurable, independent of model assumptions.
  \item Redshift falsification (Doc.~221) is \emph{strong} -- the
        $H_0$ prediction is parameter-free, several independent
        tests available.
  \item Lithium falsification (this document) is the \emph{weakest}
        -- because it builds on astrophysical models that themselves
        carry assumptions. It is more a \emph{consistency argument}
        than a sharp falsification test.
\end{itemize}

\textbf{Consequence for experimental strategy:} those wishing to
test FFGFT should begin with Casimir measurements in the
$30$--$100\,\mu$m range (Doc.~220), then with expansion-free $H_0$
measurements (GW standard sirens, Doc.~221), and only afterwards --
when the first two yield positive indications -- with cosmological
abundance consistency analyses (this document).

% ==============================================================================
\section{Summary}
% ==============================================================================

FFGFT positions itself on the cosmological lithium problem
\textbf{not} with a modified BBN reaction rate, but with a
fundamentally different view: in a static universe with
time-independent $\xi_0$, there is no BBN phase, and the observed
light element abundances result from an unbounded stellar
prehistory.

The FFGFT predictions:
\begin{itemize}[leftmargin=2em,itemsep=2pt,topsep=2pt]
  \item $\alpha$ and $\xi_0$ are constant over cosmic time.
  \item Element abundances show no systematic $z$-profile beyond
        what is explained by local chemical evolution.
  \item Pop-III stellar models reproduce $^4$He, D, $^7$Li, $^6$Li
        with a consistent set of astrophysical sources and sinks.
\end{itemize}

\textbf{Methodological note.} This position is FFGFT-internal. It
follows from the static universe (Doc.~025), the $\xi_0$ energy loss
for redshift (Doc.~064/221), and the time-independent geometric
constant (Doc.~202).

\textbf{Status.} The FFGFT position is compatible with all current
observations. It is tested by the combination of $\alpha$ variation
tests (negative, FFGFT-consistent), $z$-abundance profiles (open,
JWST-critical), and Pop-III consistency modelling (in progress,
Aanda 2025). An unambiguous refutation requires several independent
observational pillars all pointing against FFGFT -- individual
anomalies can be explained consistently.
